\begin{toreview}

\chapter{Use of Artificial Intelligence}\label{appendix:ai_usage}

In accordance with Iscte's guidelines on the use of artificial intelligence in academic work\footnote{\url{https://www.iscte-iul.pt/contents/estudantes/informacao-academica/regulamentos/2764/utilizacao-de-inteligencia-artificial}, this appendix declares the tools used, their purposes, and the author's responsibility for their outputs.}

Generative AI tools were used throughout this dissertation to support the organization of ideas, language revision, correction of typographical errors, \LaTeX{} formatting, literature-search refinement, code generation, debugging, data processing, and preliminary analysis and text for a lot of sections. Their outputs were reviewed and, where applicable, checked against the cited sources, experimental data, and executed code. The author remains fully responsible for the final text, methodology, results, references, and conclusions.

The services used were ChatGPT, Claude, Kimi, and DeepSeek. From October 2025 to September 2026, the models used through ChatGPT included GPT-5.1, GPT-5.2, GPT-5.3, GPT-5.4, GPT-5.5, GPT-5.6 Sol, GPT-6 Astra, and GPT-6 Sol. Claude models included Haiku 4.5, Opus 4.5, Opus 4.6, Sonnet 4.6, Opus 4.7, Opus 4.8, Sonnet 5, and Opus 5. Kimi models included K2 Thinking, K2.5, K2.6 and K2.7 Code. DeepSeek models included V3.2, V4, and V4.1 Flash. Free access was used for Claude, Kimi, and DeepSeek, while paid ChatGPT access was used during the final revision period. Since these services sometimes automatically selected or changed their underlying models, the exact model could not be identified for every interaction.

Limitations of these models included potential hallucinations, restricted context windows, and no persistent memory across most sessions. Their outputs were therefore treated as suggestions rather than authoritative information. Generative models tend to also amplify biases and stereotypes present in their training data \cite{bolukbasi2016man,seshadri-etal-2024-bias}. Although this risk was less prominent in the largely technical and editorial uses described here, particular attention was given to sections in which such biases could have some weight.

AI assistance occurred in every chapter, although its extent varied from spelling and formatting corrections to suggestions concerning structure, wording, code, and analysis. Substantive claims, references, numerical results, and interpretations were extensively reviewed before inclusion.

\end{toreview}